% META: {"id": "1NSI-TC-EX-025", "chapitre": "1NSI-TYPES-CONSTRUITS", "type_objet": "exercice", "capacites": ["1NSI-TYPES-CONSTRUITS-C5"], "parcours": 1, "competences": ["analyser"], "duree_min": 8, "mode_creation": "ex_nihilo", "sources_inspiration": [], "status": "approved", "capacites_codes": ["C5"], "methodes": ["M1"], "fichier_tex": "chapitres/1NSI-TYPES-CONSTRUITS/exercices/1NSI-TC-EX-025.tex", "corrige_tex": "chapitres/1NSI-TYPES-CONSTRUITS/corriges/1NSI-TC-CO-025.tex"}
% BEGIN-TRACE
% grille = [[1, 2, 3],
%            [4, 5, 6]]
% print(grille[0][1])
% print(grille[1][2])
% print(len(grille))
% print(len(grille[0]))
% EXPECTED
% 2
% 6
% 2
% 3
% END-TRACE
\begin{exercice}{1NSI-TC-EX-025}{1}{5}
On donne la grille suivante :
\begin{python}
grille = [[1, 2, 3],
           [4, 5, 6]]
print(grille[0][1])
print(grille[1][2])
print(len(grille))
print(len(grille[0]))
\end{python}
\begin{enumerate}
  \item Donner, sans exécuter, les quatre affichages.
  \item Combien de lignes et de colonnes possède cette grille ?
  \item Donner l'expression permettant d'accéder à la valeur \lstinline{4}.
\end{enumerate}
\end{exercice}
